\Needspace{10\baselineskip}
\section{Contact support in a big background}
\label{sec:big-continuous-contact}

Let $\theta$ now be an arbitrary smooth big background. Regularization makes the envelope of a continuous obstacle continuous off the poles of a strict analytic barrier. A homogeneous replacement below the obstacle then forces its Monge--Amp\`ere measure to vanish away from contact.

Choose $\rho$ with analytic singularities and smooth remainder so that
\[
 \theta+\ddc\rho\geq\kappa\omega,\qquad\kappa>0,
\]
and set $Z=\{\rho=-\infty\}$. Such a barrier exists by \cref{thm:resolvable-regularization-input}.

\Needspace{12\baselineskip}
\begin{theorem}\label{thm:big-continuous-envelope-contact}
For $h\in C(X)$, the envelope $p=P_\theta(h)$ is genuine and continuous on $X\setminus Z$. Moreover its extended Monge--Amp\`ere measure satisfies
\[
 \mathbf1_{\{p<h\}}\mathrm{MA}_\theta(p)=0.
\]
The measure in this formula is allowed to have infinite total mass.
\end{theorem}

\begin{proof}
Subtract a constant from $\rho$ so that $\rho+2\leq h$. The elementary envelope construction and upper semicontinuity of the obstacle give
\[
 \rho+2\leq p\leq h.
\]
In particular $p$ is locally bounded off $Z$. We first prove continuity there, without using the contact assertion.

Apply \cref{lem:smooth-background-regularization-rate} to $p-\sup_Xp$ and restore the constant. The transforms $\Phi_{c,\delta}$ from that lemma are defined on $X$ and, with fixed $A,K,C>0$ and $\epsilon(\delta)=C\delta$, satisfy
\begin{equation}\label{eq:big-envelope-regularization}
 \theta+\ddc\Phi_{c,\delta}\geq-(Ac+\epsilon(\delta))\omega.
\end{equation}
Their infimum is attained at a positive radius off $Z$, and
\[
 p-K(\delta^2+\delta)\leq\Phi_{c,\delta}
       \leq R_\delta p\leq R_\delta h.
\]
For small $\delta$ the first bound gives $\Phi_{c,\delta}\geq\rho+1$, hence $\rho\leq\Phi_{c,\delta}$ on $X$.

\Needspace{5\baselineskip}
Set
\[
 m(\delta)=\sup_{0<t\leq\delta}\|R_th-h\|_\infty,
 \qquad c=m(\delta)+\epsilon(\delta)+\delta,
 \qquad s=\frac{Ac+\epsilon(\delta)}{\kappa}.
\]
Positivity, normalization, and localization of the averaging kernels imply $m(\delta)\to0$ by uniform continuity of $h$. Define
\[
 v_\delta=\frac{\Phi_{c,\delta}+s\rho}{1+s}.
\]
The strict barrier and \eqref{eq:big-envelope-regularization} show that $v_\delta$ is a genuine $\theta$-plurisubharmonic function. Since $\rho\leq\Phi_{c,\delta}$,
\[
 v_\delta\leq\Phi_{c,\delta}\leq R_\delta h\leq h+m(\delta).
\]
Thus $v_\delta-m(\delta)$ is an envelope competitor, and $v_\delta\leq p+m(\delta)$.

On a fixed compact subset $L\Subset X\setminus Z$, the difference $p-\rho$ is bounded above. As $s\leq Cc$, the exact identity
\[
 \Phi_{c,\delta}-p=(1+s)(v_\delta-p)+s(p-\rho)
\]
therefore gives $\Phi_{c,\delta}-p\leq C_Lc$. At a minimizing radius $t$,
\[
 -c\log(t/\delta)
        \leq\Phi_{c,\delta}-p+K(\delta^2+\delta)\leq C_Lc.
\]
Since $c\geq\delta$, this yields $t\geq a_L\delta$ with $a_L>0$ fixed. Monotonicity of smoothing now implies
\[
 p-K(a_L\delta)^2\leq R_{a_L\delta}p
       \leq\Phi_{c,\delta}+K\delta^2+K\delta.
\]
Consequently $R_{a_L\delta}p\to p$ uniformly on $L$. These averages are smooth for positive radius even though $p$ is globally unbounded. Applying this to compact neighborhoods proves continuity on $X\setminus Z$.

Now take $x\in\{p<h\}\setminus Z$. By the continuity just proved this is an ordinary open set near $x$. Choose a sufficiently small coordinate ball $B_r$ centered at $x$, with closure in that set. Let $H$ be the homogeneous replacement of \cref{thm:ball-homogeneous-replacement} with boundary data $p|_{\partial B_r}$. Its domination property gives $H\geq p$ on the ball.

We can make the ball small enough that $H\leq h$. Indeed, with $z$ denoting the chosen Euclidean coordinates and $\Delta$ their scalar Laplacian, $\theta+\ddc H\geq0$ implies a fixed bound $\Delta H\geq-C$. Comparison with the function
$\sup_{\partial B_r}p+C'(r^2-|z-x|^2)$, with $C'$ fixed large enough, gives
\[
 H\leq\sup_{\partial B_r}p+C'r^2.
\]
As $r\to0$, the right-hand side tends to $p(x)$, whereas $\inf_{B_r}h\to h(x)>p(x)$. This proves the claim for sufficiently small $r$.

Paste $H$ into $p$ on the ball. The result is $\theta$-plurisubharmonic: for each $a>0$, paste $\max\{p,H-a\}$ inside the ball and $p$ outside. Continuity and equality of the boundary values ensure that $H-a<p$ on a collar, so this is an ordinary local maximum construction. Letting $a\downarrow0$ gives the desired pasted function, with convergence locally uniform near the boundary. It is bounded above by $h$ on all of $X$. The defining property of the envelope therefore forces it to equal $p$. Thus $p=H$ on $B_r$ and $\mathrm{MA}_\theta(p)$ vanishes there.

These balls cover $\{p<h\}\setminus Z$. The local bounded product there agrees with the extended measure, and the latter does not charge the analytic, hence pluripolar, set $Z$. This proves the contact assertion on $X$, without taking or bounding its total mass.
\end{proof}
